% Created 2026-10-27 Tue 15:32
% Intended LaTeX compiler: pdflatex
\documentclass[11pt]{article}
\usepackage[utf8]{inputenc}
\usepackage[T1]{fontenc}
\usepackage{graphicx}
\usepackage{grffile}
\usepackage{longtable}
\usepackage{wrapfig}
\usepackage{rotating}
\usepackage[normalem]{ulem}
\usepackage{amsmath}
\usepackage{textcomp}
\usepackage{amssymb}
\usepackage{capt-of}
\usepackage{hyperref}
\author{Arnab Chakraborty}
\date{\today}
\title{}
\hypersetup{
 pdfauthor={Arnab Chakraborty},
 pdftitle={},
 pdfkeywords={},
 pdfsubject={},
 pdfcreator={Emacs 26.3 (Org mode 9.1.9)}, 
 pdflang={English}}
\begin{document}

\tableofcontents

\section{Step by step creation of an Android app}
\label{sec:org30812e2}

\subsection{List of tools:}
\label{sec:org5979edc}
\begin{center}
\begin{tabular}{lll}
Tools & Source & Doc\\
\hline
javac, jar, keytool & JDK & \\
android.jar & \$ANDROID\(_{\text{HOME}}\)/platforms/android-35/ & \url{https://developer.android.com/reference}\\
aapt2 & \$ANDROID\(_{\text{HOME}}\)/build-tools/34.0.0/ & \url{https://developer.android.com/tools/aapt2}\\
d8 & ditto & \url{https://developer.android.com/tools/d8}\\
zipalign & ditto & \url{https://developer.android.com/tools/zipalign}\\
apksigner & ditto & \url{https://developer.android.com/tools/apksigner}\\
adb & \$ANDROID\(_{\text{HOME}}\)/platform-tools/ & \url{https://developer.android.com/tools/adb}\\
 &  & \\
\end{tabular}
\end{center}


\subsection{Check what you already have:}
\label{sec:orgaf3f036}
echo "\$ANDROID\(_{\text{HOME}}\)"

ls "\$ANDROID\(_{\text{HOME}}\)/platforms/android-35/"
ls "\$ANDROID\(_{\text{HOME}}\)/build-tools/34.0.0/"

javac -version

\subsection{If needed}
\label{sec:orgc38eb1f}
sdkmanager "build-tools;34.0.0"
sdkmanager "platform-tools"

\subsection{A little note}
\label{sec:org6391ec8}
One small but important detail: android.jar is not a complete
Android runtime. It supplies the API classes needed to compile
your application against Android. Your phone supplies the actual
runtime when the app executes. 

\section{The process starts}
\label{sec:org7d8a4be}

\subsection{Stage 1A — Java compilation and inspecting a .class file.}
\label{sec:org5c834dd}

Aim: Compile a tiny Java class in the usual way. Inspect the
generated .class file using javap.

Create a new folder and write java code:

\begin{verbatim}
public class Hello {
    public static void main(String[] args) {
        System.out.println("Hello, Android!");
    }
}
\end{verbatim}

Compile as usual:

\begin{verbatim}
javac Hello
\end{verbatim}

Inspect the .class file:

\begin{verbatim}
javap -c Hello
javap -verbose Hello
\end{verbatim}

\subsection{Stage 1B — DEX}
\label{sec:org60e8fbe}

Use d8 to turn the Java .class files into classes.dex. Inspect
the result. 

\begin{verbatim}
mkdir -p dex

"$ANDROID_HOME/build-tools/34.0.0/d8" \
    --output dex \
    Hello.class
\end{verbatim}

Inspect it:
\begin{verbatim}
"$ANDROID_HOME/build-tools/34.0.0/dexdump" \
                    -d dex/classes.dex
\end{verbatim}
This requires a tool dexdump not listed in the table
earlier. This inspection step is optional.

\subsection{Stage 2A - Compiling against android.jar}
\label{sec:org08178ef}
\begin{verbatim}
mkdir -p ~/campapk/stage2/src/com/example/helloworld
cd ~/campapk/stage2
\end{verbatim}

Write the following java code in the helloworld folder:
\begin{verbatim}
package com.example.helloworld;

import android.app.Activity;
import android.os.Bundle;
import android.widget.TextView;

public class MainActivity extends Activity {
    @Override
    protected void onCreate(Bundle savedInstanceState) {
        super.onCreate(savedInstanceState);

        TextView text = new TextView(this);
        text.setText("Hello World!");
        text.setTextSize(32);

        setContentView(text);
    }
}
\end{verbatim}

From the stage2 folder run:
\begin{verbatim}
mkdir -p classes

javac \
  -classpath "$ANDROID_HOME/platforms/android-35/android.jar" \
  -d classes \
  src/com/example/helloworld/MainActivity.java
\end{verbatim}
\subsection{Stage 2B — Android resources and manifest}
\label{sec:org50ad3bc}

An ordinary Java main() method is not the entry point Android
uses to launch an application. Android launches components
declared in the manifest, such as an Activity 

Write AndroidManifest.xml, use aapt2 to compile and link
resources, and examine the generated R.java and binary manifest. 

Our next job is to tell Android what our application is, which
Activity it can launch, and how its resources are packaged. 

Create the following AndroidManifest.xml file in the stage2
folder:
\begin{verbatim}
<?xml version="1.0" encoding="utf-8"?>
<manifest xmlns:android="http://schemas.android.com/apk/res/android"
    package="com.example.helloworld">

    <application
        android:label="Hello World"
        android:theme="@android:style/Theme.Material.Light">

        <activity
            android:name="com.example.helloworld.MainActivity"
            android:exported="true">

            <intent-filter>
                <action android:name="android.intent.action.MAIN" />
                <category android:name="android.intent.category.LAUNCHER" />
            </intent-filter>

        </activity>
    </application>
</manifest>
\end{verbatim}

Explanation: 
A few important details:

\begin{itemize}
\item android:label gives the application its displayed name.

\item android:name identifies the Activity class we just compiled.

\item MAIN and LAUNCHER tell Android that this is the Activity to launch from the app launcher.

\item android:exported="true" is required for an Activity with this
\end{itemize}
intent filter when targeting modern Android versions. 

\subsubsection{{\bfseries\sffamily TODO} The last two points are not clear to me.}
\label{sec:orgfd29631}

Create two folders:
\begin{verbatim}
mkdir -p res-compiled linked
\end{verbatim}

Now issue the main command:
\begin{verbatim}
"$ANDROID_HOME/build-tools/34.0.0/aapt2" link \
  -o linked/base.apk \
  --manifest AndroidManifest.xml \
  -I "$ANDROID_HOME/platforms/android-35/android.jar" \
  --min-sdk-version 23 \
  --target-sdk-version 35
\end{verbatim}

Produces base.apk in the linked folder. Let's inspect it:
\begin{verbatim}
unzip -l linked/base.apk
\end{verbatim}
The output is
\begin{verbatim}
Archive:  linked/base.apk
  Length      Date    Time    Name
---------  ---------- -----   ----
     1788  1980-01-01 05:30   AndroidManifest.xml
       40  1980-01-01 05:30   resources.arsc
---------                     -------
     1828                     2 files

\end{verbatim}
This still does not have the MainActivity bytecode. We shall
first turn that into hex.
\begin{verbatim}
mkdir -p dex

"$ANDROID_HOME/build-tools/34.0.0/d8" \
  --lib "$ANDROID_HOME/platforms/android-35/android.jar" \
  --min-api 23 \
  --output dex \
  classes/com/example/helloworld/MainActivity.class
\end{verbatim}
This produces classes.dex in the dex folder.

\subsubsection{Why these options?}
\label{sec:org41cc5c9}
\begin{itemize}
\item --lib android.jar tells D8 where Android's API classes are
defined. They are provided by the Android system, so we don't
bundle them into our app.
\item --min-api 23 sets the minimum Android API level for the generated bytecode.
\item --output dex puts the resulting DEX file in the dex directory.
\end{itemize}

\subsection{Stage 2C — Assemble the APK}
\label{sec:org27ad749}

Combine the compiled resources, manifest, DEX file and other
required files into an APK. Understand the APK's ZIP structure.  

We shall now put classes.dex inside the APK
An APK is essentially a ZIP archive with particular required
contents. We can add our DEX file to the APK that aapt2 created. 

\begin{verbatim}
cp linked/base.apk app-unsigned.apk

zip -j app-unsigned.apk dex/classes.dex
\end{verbatim}


\subsection{Stage 2D — Align}
\label{sec:org286de6c}

Run zipalign, create a development key with keytool, sign with
apksigner, and verify the signature. 

Android expects APK contents to be aligned appropriately for
efficient access. The SDK provides zipalign for this. 

\begin{verbatim}
"$ANDROID_HOME/build-tools/34.0.0/zipalign" \
  -f -v 4 \
  app-unsigned.apk \
  app-aligned.apk
\end{verbatim}

Let's check the alignment:

\begin{verbatim}
"$ANDROID_HOME/build-tools/34.0.0/zipalign" \
  -c -v 4 app-aligned.apk
\end{verbatim}

\subsection{Stage 2E -- Sign}
\label{sec:org1c42c1e}
\begin{verbatim}
keytool -genkeypair \
  -keystore debug.keystore \
  -alias androiddebugkey \
  -keyalg RSA \
  -keysize 2048 \
  -validity 10000 \
  -storepass android \
  -keypass android \
  -dname "CN=Android Debug,O=Android,C=US"
\end{verbatim}

Now sign:

\begin{verbatim}
"$ANDROID_HOME/build-tools/34.0.0/apksigner" sign \
  --ks debug.keystore \
  --ks-key-alias androiddebugkey \
  --ks-pass pass:android \
  --key-pass pass:android \
  --out app-debug.apk \
  app-aligned.apk
\end{verbatim}

Check the signature:
\begin{verbatim}
"$ANDROID_HOME/build-tools/34.0.0/apksigner" verify \
  --verbose \
  --print-certs \
  app-debug.apk
\end{verbatim}
\subsection{Stage 6 — Install and run}
\label{sec:org37fef17}

Use adb install to put the APK on your Android 16 phone and launch it.
\end{document}
